% Shared preamble for every document under content/.
%
% Documents are compiled twice:
%   * tex4ht (make4ht) turns them into HTML fragments for the frontend shell
%   * pdflatex still produces a normal PDF, so every helper below degrades
%     to a no-op when LaTeX is not running under tex4ht.
%
% The frontend shell owns layout, navigation and the terminal aesthetic.
% Everything below the "content vocabulary" heading is presentation that
% LaTeX cannot express on its own (videos, pills, cards).

\documentclass[11pt]{article}

\usepackage[T1]{fontenc}
\usepackage[utf8]{inputenc}
\usepackage{lmodern}
\usepackage{graphicx}
\usepackage{amsmath}
\usepackage{amssymb}
\usepackage{longtable}
\usepackage[hidelinks]{hyperref}

% The site renders plain section titles, never numbered ones.
\setcounter{secnumdepth}{0}

% ---------------------------------------------------------------------------
% tex4ht plumbing
% ---------------------------------------------------------------------------

% tex4ht defines \HCode; plain LaTeX does not, so provide a no-op fallback.
\newif\ifhtmlmode
\ifdefined\HCode\htmlmodetrue\else\htmlmodefalse\fi
\providecommand{\HCode}[1]{}

% Emit raw HTML. No-op in PDF mode.
\newcommand{\hcode}[1]{\ifhtmlmode\HCode{#1}\fi}

% Declares the language of the code block that follows. The HTML converter
% reads the marker and highlights the block with Pygments.
\newcommand{\codelang}[1]{\hcode{<!--lang:#1-->}}

% ---------------------------------------------------------------------------
% Content vocabulary
% ---------------------------------------------------------------------------

% A tag pill, e.g. \pill{Go}
%
% The text is typeset outside the raw HTML so that LaTeX escapes apply to it.
% tex4ht folds adjacent identical <span> elements into one, so the pill uses
% an <i> element and carries the "pill" class the stylesheet keys off.
\newcommand{\pill}[1]{\hcode{<i class="pill">}#1\hcode{</i>}\hcode{ }}

% A row of pills. The pills are passed in already rendered, so that no
% \hcode call is ever nested inside another one.
\newcommand{\pillrow}[1]{\hcode{<div class="pill-row">}#1\hcode{</div>}}

% An image served from the site's static assets.
%   \webimg{src}{alt}{class}
%   \webimgw{src}{alt}{class}{width}
\newcommand{\webimg}[3]{\hcode{<img src="#1" alt="#2" class="#3" loading="lazy">}}
\newcommand{\webimgw}[4]{\hcode{<img src="#1" alt="#2" class="#3" width="#4" loading="lazy">}}

% An autoplaying, looping video served from the site's static assets.
\newcommand{\webvideo}[2]{\hcode{<video class="project-video" src="#1" autoplay muted loop playsinline tabindex="-1" preload="metadata" aria-label="#2"></video>}}

% A link that opens in a new tab.
\newcommand{\extlink}[2]{\hcode{<a href="#1" target="_blank" rel="noopener noreferrer"><strong>}#2\hcode{</strong></a>}}

% A company logo next to an employer name.
\newcommand{\orglogo}[2]{\hcode{<img class="org-logo org-logo--#1" src="/assets/#2" alt="#1">}}

% The date line of an experience entry.
\newcommand{\orgdate}[1]{\hcode{<br/><span class="date">#1</span>}}

% --- open source contributions ---------------------------------------------

% The title of an open source contribution, linking to the pull request.
\newcommand{\osslink}[2]{\hcode{<a href="#2" target="_blank" rel="noopener noreferrer"><strong>}#1\hcode{</strong></a><br/>}}

\newcommand{\ossdesc}[1]{\hcode{<span>}#1\hcode{</span>}}

% --- project cards ---------------------------------------------------------

\newenvironment{projectcard}{%
  \hcode{<article class="project-card">}%
}{%
  \hcode{</article>}%
}

\newcommand{\projecttitle}[2]{%
  \hcode{<div class="project-body"><a href="#2" target="_blank" rel="noopener noreferrer"><strong>}#1\hcode{</strong></a>}%
}

\newcommand{\projectdesc}[1]{\hcode{<p>}#1\hcode{</p>}}

% A link that stays on the site, usable inside \projectdesc.
\newcommand{\sitelink}[2]{\hcode{<a href="#1">}#2\hcode{</a>}}

% A literal underscore. LaTeX's \_ does not survive the HTML conversion.
\newcommand{\uscore}{\char95}

% A literal tilde. LaTeX's \textasciitilde does not survive the conversion.
\newcommand{\tildechar}{\char126}

% A horizontal rule.
\newcommand{\webhr}{\hcode{<hr/>}}

% --- inline links in prose -------------------------------------------------

% A link that stays on the site. Use \href inside ordinary prose.
